\documentclass[10.5pt,letterpaper]{article}
\usepackage[utf8]{inputenc}
\usepackage[T1]{fontenc}
\usepackage{geometry}
\geometry{margin=0.82in,top=0.72in,bottom=0.72in}
\usepackage{hyperref}
\usepackage{setspace}
\usepackage{siunitx}
\DeclareSIUnit\calorie{cal}
\sisetup{detect-all=true, separate-uncertainty=true, range-phrase=--}
\setstretch{1.03}
\setlength{\parindent}{0pt}
\setlength{\parskip}{0.48em}

\begin{document}

\textbf{Myke Vital Sao Temgoua}\\
Department of Physics, Faculty of Science\\
University of Yaound\'{e} I, P.O. Box 812, Yaound\'{e}, Cameroon\\
\href{mailto:myke-vital.sao@facsciences-uy1.cm}{myke-vital.sao@facsciences-uy1.cm}
\hfill August 9, 2026

\vspace{0.35em}
\textbf{Prof. Kenneth M. Merz, Jr., Editor-in-Chief}\\
\textit{Journal of Chemical Information and Modeling}\\
American Chemical Society

\vspace{0.35em}
Dear Prof. Merz,

We submit our manuscript, \textbf{``Computational Discovery of Antimalarial Candidates from African Natural Product Chemical Space''}, for consideration as an Article in \textit{JCIM}. The manuscript presents a modelling-only workflow for resource-efficient antimalarial discovery that combines VAE-guided chemical-space exploration, centroid-based docking, target-anchored multi-target modelling, and resistance-aware computational prioritization.

The first contribution is methodological efficiency. A \num{65856}-molecule hybrid library derived from \num{396} African natural products and \num{454} synthetic antimalarials was represented by \num{484} chemically diverse centroids, reducing the nominal number of docking representatives by \qty{99.3}{\percent}. The resulting analysis quantifies whole-molecule novelty (\qty{92.6}{\percent} of the evaluated sample below a maximum seed Tanimoto similarity of \num{0.4}) together with \qty{69.3}{\percent} recovery of seed scaffolds. These results are presented as computational chemical-space measurements, not as evidence of biological activity.

The second contribution is an explicit target-anchored polypharmacology layer. A predefined cohort of \num{17} polypharmacology candidates was evaluated against PfDHFR, PfCRT, PfClpP, and PfATP4 using corrected target-specific anchors and fixed Vina boxes. The \num{68} candidate--target computations passed the declared pose-and-anchor gate, yielding target-wise mean docking scores of \qty{-5.755}{\kilo\calorie\per\mole} (PfDHFR), \qty{-6.185}{\kilo\calorie\per\mole} (PfCRT), \qty{-5.935}{\kilo\calorie\per\mole} (PfClpP), and \qty{-6.308}{\kilo\calorie\per\mole} (PfATP4). These values are scoring-function outputs, not experimental binding free energies.

The third contribution is resistance-aware computational prioritization. Using a corrected per-target docking-derived RRS definition that excludes non-binding WT baselines, the \num{17}-candidate cohort spans four RRS classes (A*: \num{6}, B: \num{5}, C: \num{5}, D: \num{1}). The same cohort was characterized by PNS and ACSI descriptors, with PNS ranging from \num{1.04} to \num{6.00} and ACSI from \num{0.173} to \num{0.823}. We interpret these quantities as hypothesis-generating computational descriptors: they do not establish mutation resilience, simultaneous target engagement, therapeutic benefit, or experimental efficacy.

The manuscript is deliberately transparent about model boundaries. The corrected PfClpP structure is PDB 2F6I rather than the historically misassigned 4GM2/PfClpR entry; PfCRT uses a Y01 membrane-proxy cavity, and PfATP4 uses a motif-defined catalytic-region anchor because PDB 9N10 lacks a co-crystallized small-molecule inhibitor. Legacy DiffDock consensus artifacts and the unsigned consensus run is excluded. All source files, hashes, scripts, raw Vina outputs, the integrated Vina--RRS--PNS--ACSI table, and its reproducibility metadata are available in the repository. No experimental activity, affinity, MD-RRS, or biological validation claim is made.

We believe the work is appropriate for \textit{JCIM} because it contributes a reproducible, resource-aware modelling framework that links generative chemical-space exploration to target-anchored multi-target and resistance-aware prioritization. The manuscript is not intended to substitute computational scores for experiment; rather, it defines a transparent hypothesis space and auditable candidate-selection workflow for subsequent medicinal-chemistry and biological studies.

This manuscript represents original work, is not under consideration elsewhere, and has been reviewed and approved by all authors. The authors declare no competing interests.

\vspace{0.35em}
Sincerely,\\[0.25em]
\textbf{Myke Vital Sao Temgoua}\\
(on behalf of all co-authors)

\end{document}
